% Inhaltsverzeichnis
\newpage
\pagestyle{empty}
\renewcommand{\contentsname}{} % Entfernt die Standardüberschrift
\phantomsection
\section*{Inhaltsverzeichnis}
\addcontentsline{toc}{section}{Inhaltsverzeichnis} % Fügt zum Inhaltsverzeichnis hinzu
\vspace{-3cm} % Abstand verringern
\tableofcontents
\newpage

% Abbildungsverzeichnis
\newpage
\pagestyle{empty}
\renewcommand{\listfigurename}{} % Entfernt die Standardüberschrift
\phantomsection
\section*{Abbildungsverzeichnis}
\addcontentsline{toc}{section}{Abbildungsverzeichnis} % Fügt zum Inhaltsverzeichnis hinzu
\vspace{-1.5cm} % Abstand verringern
\listoffigures
\newpage

% Tabellenverzeichnis
\newpage
\pagestyle{empty}
\renewcommand{\listtablename}{} % Entfernt die Standardüberschrift
\phantomsection
\section*{Tabellenverzeichnis}
\addcontentsline{toc}{section}{Tabellenverzeichnis}
\vspace{-1.5cm} % Abstand verringern
\listoftables
\newpage

% Abkürzungsverzeichnis (keine Kopfzeile, keine Linie)
\newpage
\pagestyle{fancy} % Hier wieder auf fancy stellen, um Seitenzahlen zu zeigen
\fancyhead{}
\renewcommand{\headrulewidth}{0pt}  % Entfernt die Linie in der Kopfzeile
\phantomsection
\section*{Abkürzungsverzeichnis}
\addcontentsline{toc}{section}{Abkürzungsverzeichnis}
Hier steht das Abkürzungsverzeichnis.

% Symbolverzeichnis (keine Kopfzeile, keine Linie)
\newpage
\phantomsection
\section*{Symbolverzeichnis}
\addcontentsline{toc}{section}{Symbolverzeichnis}
Hier steht das Symbolverzeichnis.

% Formelverzeichnis (keine Kopfzeile, keine Linie)
\newpage
\phantomsection
\section*{Formelverzeichnis}
\addcontentsline{toc}{section}{Formelverzeichnis}
Hier steht das Formelverzeichnis.
